\subsection{Projective Non-Termination and the Condition of Temporal Ordering}
  \label{subsec:projective-non-termination}

  Time is identified with the ordering of distinguishable projected
  states: $\Delta \tau \sim \mathrm{dist}(U_n, U_{n+1})$.
  A hypothetical terminal state~$U_0$ satisfying
  $\Pi(\chi + \delta\chi) = U_0$ for all admissible~$\delta\chi$ would
  imply $\Delta\tau \rightarrow 0$ for all internal observers.

  The inaccessibility of absolute zero emerges as a projective necessity:
  zero-point fluctuations are reinterpreted as the minimal projective
  bandwidth required to prevent collapse of the observable description into
  a stationary state.
  Time ceases only when the projection becomes stationary---a regime
  structurally forbidden by the existence of a finite bound on projective
  resolvability.

  The connection between projective stationarity and injectivity can now be made
  explicit.
  A stationary projection corresponds to the limit in which, within any admissible
  configuration neighborhood, distinct substrate configurations fail to produce
  distinguishable projected states: no local variation $\delta\chi$ induces a resolvable
  change in the effective description, so that the accessible neighborhood is mapped into an effectively
  indistinguishable projected class and no update can be resolved.
  This is the opposite extreme from injectivity, but leads to the same operational
  consequence: $U_{n+1} = U_n$, since neither extreme admits a non-trivial ordering of
  distinguishable projected states.
  This equivalence unifies three results established independently in this paper:
  the inaccessibility of absolute zero (Section~\ref{subsec:absolute-zero-necessity}), the suppression of vacuum
  fluctuations in injective regimes (Section~9.4), and the projective non-termination
  condition stated here.
  All three express the same structural fact: $S_\Pi > 0$ is a necessary condition
  for temporal ordering, and $S_\Pi = 0$ is the universal limit at which time, energy,
  and fluctuations simultaneously vanish.
  Locally, this limit can be approached---corresponding to quasi-classical or
  decohered regimes---but it cannot be globally reached without eliminating the
  effective description itself.
